% pustaka.tex — dari DAFTAR_PUSTAKA.md (tanpa heading), hanging indent
\sloppy
\raggedright

\begin{enumerate}[leftmargin=1.5cm,itemindent=-1.5cm,label={},itemsep=4pt,parsep=2pt]
    \item Abbas, S. F., \& Rawabdeh, M. A. (2025). The concept of \textit{munkar} in al-Dhahabi's critique of al-Hakim's hadith authentication in al-\textit{Mustadrak}. \textit{Jurnal Studi Ilmu-Ilmu Al-Qur'an dan Hadis}, \textit{26}(2), 545--568. \url{https://doi.org/10.14421/qh.v26i2.6267}
    \item Abdullah, Z., Shafira, N., Ibrahim, I., \& Rosdiana, R. (2026). Viralitas dan validitas: Digitalisasi \textit{takhrij} hadis di era media sosial. \textit{Dirayah: Jurnal Ilmu Hadis}, \textit{6}(2), 381--390. \url{https://doi.org/10.62359/dirayah.v6i02.1335}
    \item Abdulrahman, M. A. (2024). The future of hadith studies in the digital age: Opportunities and challenges. \textit{Journal of Ecohumanism}, \textit{3}(8), [Halaman --- VERIFIKASI\_MANUAL]. \url{https://doi.org/10.62754/joe.v3i8.4927}
    \item Al-Imam, A. A. R. (2026). The influence of modernity on \textit{matn} criticism: A comparative study of traditionist and reformist approaches in Ibn Hajar's \textit{Fath al-Bari} and Rashid Rida's \textit{Majallat al-Manar} on hadiths of natural phenomena. \textit{Journal of Islamic Thought and Civilization}, \textit{16}(1), 01--17. \url{https://doi.org/10.32350/jitc.161.01}
    \item Alam, T., \& Schneider, J. (2020). Social network analysis of hadith narrators from \textit{Sahih Bukhari}. \textit{Proceedings of 2020 7th IEEE International Conference on Behavioural and Social Computing (BESC)}, 1--6. \url{https://doi.org/10.1109/BESC51023.2020.9348299}
    \item Alghamdi, J., Albukhari, A., \& Al-Dala'in, T. (2025). Pretrained models against traditional machine learning for detecting fake hadith. \textit{Electronics}, \textit{14}(17), 3484. \url{https://doi.org/10.3390/electronics14173484}
    \item Ardig, N., Karagezoglu, M. M., et al. (2025). Genre analysis and religious texts: A methodological model of hadith commentary. \textit{Journal of Islamic Studies}, \textit{36}(2), 163--218. \url{https://doi.org/10.1093/jis/etae068}
    \item Azmi, A. M., Al-Qabbany, A. O., \& Hussain, A. (2019). Computational and natural language processing based studies of hadith literature: A survey. \textit{Artificial Intelligence Review}, \textit{52}(2), 1369--1414. \url{https://doi.org/10.1007/s10462-019-09692-w}
    \item Bin Jamil, K. H. (2024). Between traditionalising and futuring: Applying the broader \textit{maqasid} paradigm to hadith studies. \textit{Asia Pacific Journal of Educators and Education}, \textit{39}(2), 183--196. \url{https://doi.org/10.21315/apjee2024.39.2.10}
    \item Calgan, M. A. (2024). Analysis of Ibn Kathir's content criticism of some of the \textit{Sahihayn} hadiths in the \textit{sirah} section of \textit{al-Bidaya wa al-Nihaya}. \textit{Cumhuriyet Ilahiyat Dergisi}, \textit{28}(1), 303--324. \url{https://doi.org/10.18505/cuid.1430992}
    \item Gil, R. E. (2024). Tracing the \textit{rihla} in \textit{buldan} books: Ibn al-Faqih's \textit{muhaddith} resources. \textit{Hitit Theology Journal}, \textit{23}(Din ve Coğrafya), 140--154. \url{https://doi.org/10.14395/hid.1417152}
    \item Haque, F., Orthy, A. H., \& Siddique, S. (2020). Hadith authenticity prediction using sentiment analysis and machine learning. \textit{2020 IEEE 14th International Conference on Application of Information and Communication Technologies (AICT)}, 1--6. \url{https://doi.org/10.1109/aict50176.2020.9368569}
    \item Jaiyeoba, H. B., \& Osmani, N. M. (2024). Hadith preservation: Techniques and contemporary efforts. \textit{Journal of Fatwa Management and Research}, \textit{29}(3), 31--45. \url{https://doi.org/10.33102/jfatwa.vol29no3.597}
    \item Kamran, A. B., Abro, B., \& Basharat, A. (2023). SemanticHadith: An ontology-driven knowledge graph for the hadith corpus. \textit{Journal of Web Semantics}, \textit{78}, 100797. \url{https://doi.org/10.1016/j.websem.2023.100797}
    \item Kara, S. (2026). Debating the origins: The sanctity of Madina in hadith narratives. \textit{Journal of Islamic Studies}, \textit{37}(2), 163--207. \url{https://doi.org/10.1093/jis/etaf047}
    \item Liew, H. H. (2025). "The caliphate will last for thirty years": Polemic and political thought in the afterlife of a prophetic hadith. \textit{Journal of Islamic Studies}, \textit{36}(1), 38--82. \url{https://doi.org/10.1093/jis/etae049}
    \item Lutfianto, M., Abdurrahman, A., \& Takwallo, T. (2024). History of hadith writing in the precodification period in the prospective of 'Ajjaj al-Khatib. \textit{Al-Thiqah: Jurnal Ilmu Keislaman}, \textit{7}(2), 308. \url{https://doi.org/10.56594/althiqah.v7i2.263}
    \item Matin Bin Salman, A. M. B. (2025). Reconstructing hadith discourse in the digital age: From text to discourse. \textit{Journal of Ecohumanism}, \textit{4}(1), 1--11. \url{https://doi.org/10.62754/joe.v4i1.4084}
    \item Mohamed, E., \& Sarwar, R. (2022). Linguistic features evaluation for hadith authenticity through automatic machine learning. \textit{Digital Scholarship in the Humanities}, \textit{37}(3), 830--843. \url{https://doi.org/10.1093/llc/fqab092}
    \item Mosa, M. A. (2025). Synergizing structure and semantics: A knowledge graph-transformer framework for narrator disambiguation in hadith networks. \textit{Digital Scholarship in the Humanities}, [VERIFIKASI\_MANUAL --- Crossref: 404]. \url{https://doi.org/10.1093/llc/fqaf088}
    \item Muhamed Ali, M. A. M., Ibrahim, M. N., Usman, A. H., \ldots{}, Kadir, M. N. A. (2015). \textit{Al-jarh wa al-ta'dil} (criticism and praise): Its significance in the science of hadith. \textit{Mediterranean Journal of Social Sciences}, \textit{6}(2S1), 284. \url{https://doi.org/10.5901/mjss.2015.v6n2s1p284}
    \item Musadad, A. (2026). Articulating the prophetic authority: Some notes on the nature of hadith collection in \textit{al-Muwatta'} and \textit{musannaf} literature. \textit{Jurnal Studi Ilmu-Ilmu Al-Qur'an dan Hadis}, \textit{27}(1), 351--373. \url{https://doi.org/10.14421/qh.v27i1.7472}
    \item Noor, U. M. (2023). \textit{Muqaddimah} of Ibn al-Salah and the revival of hadith studies in the Mamluk era. \textit{Al-Bayan: Journal of Qur'an and Hadith Studies}, \textit{21}(2), 271--297. \url{https://doi.org/10.1163/22321969-20230135}
    \item Obeid, S. M. S., \& Hussein, F. S. (2023). Al-Bukhari's sources. \textit{Jordan Journal for History and Archaeology}, \textit{17}(3), 44--62. \url{https://doi.org/10.35516/jjha.v17i3.434}
    \item Rashwan, H. (2024). Hadith as oral literature through early Islamic literary criticism. \textit{Studia Islamica}, \textit{119}(1), 34--69. \url{https://doi.org/10.1163/19585705-12341481}
    \item Razak, N. M. A. M., Nazri, M. A., Latifah, L. A., \ldots{}, Ibrahim, N. S. (2025). Evolution of hadith textual criticism: From classical to contemporary approaches. \textit{Global Journal Al-Thaqafah}, 15--29. \url{https://doi.org/10.7187/GJATSI122025-2}
    \item Rozikin, M. R. (2026). \textit{Al-fara'id} is half of knowledge: A critical hadith analysis and its juridical implications in Islamic inheritance discourse. \textit{Cogent Arts and Humanities}, \textit{13}(1), [Halaman --- VERIFIKASI\_MANUAL]. \url{https://doi.org/10.1080/23311983.2025.2598289}
    \item Shaaban, M., Elshenawy, A., \& El-Din, S. (2026). Prophetic authorial style modeling for detecting fabricated hadiths using AraBERT. \textit{Journal of Computing \& Biomedical Informatics}, \textit{10}(2), [Halaman --- VERIFIKASI\_MANUAL]. \url{https://doi.org/10.56979/1002/2026/1315}
    \item Su, I.-W. (2021). 'Ali b. al-Madini (161--234/778--849): A critical reconstruction of his biography and evaluation of his contribution to hadith criticism. \textit{Journal of Islamic Studies}, \textit{33}(1), 1--34. \url{https://doi.org/10.1093/jis/etab049}
    \item Supriyadi, T., Julia, J., Aeni, A. N., \& Sumarna, E. (2020). Action research in hadith literacy: A reflection of hadith learning in the digital age. \textit{International Journal of Learning, Teaching and Educational Research}, \textit{19}(5), 99--124. \url{https://doi.org/10.26803/ijlter.19.5.6}
    \item Syukur, H., Nasrulloh, N., \& Fachmuroji, I. (2026). Living hadith in the age of artificial intelligence: A phenomenological study of the use of ChatGPT as a source for studying hadith among college students. \textit{Asian Journal of Islamic Studies and Da'wah}, \textit{4}(3), 189--204. \url{https://doi.org/10.58578/ajisd.v4i3.10228}
    \item Tanjung, T. N. S., \& Fadhlurrahman, Z. (2025). The role of \textit{tadwin} hadith in the development of Islamic historiography. \textit{Istifham: Journal of Islamic Studies}, \textit{3}(3), 188--201. \url{https://doi.org/10.71039/istifham.v3i3.123}
    \item Tarmom, T., Atwell, E., \& Alsalka, M. (2022). Deep learning vs compression-based vs traditional machine learning classifiers to detect hadith authenticity (pp. 206--222). \textit{Communications in Computer and Information Science}. \url{https://doi.org/10.1007/978-3-031-04447-2_14}
    \item Wasman, Mesraini, M., \& Suwendi, S. (2023). A critical approach to prophetic traditions: Contextual criticism in understanding hadith. \textit{Al-Jami'ah: Journal of Islamic Studies}, \textit{61}(1), 1--17. \url{https://doi.org/10.14421/ajis.2023.611.1-17}
    \item Watukila, S. (2024). Historiografi hadis di masa \textit{mutaqaddimin} dalam pandangan Mustafa Azami. \textit{Amsal Al-Qur'an: Jurnal Al-Qur'an dan Hadis}, \textit{1}(3), 221--235. \url{https://doi.org/10.63424/amsal.v1i3.126}
\end{enumerate}
